\documentclass[10pt]{article}
\usepackage[a4paper,margin=0.6in]{geometry}
\usepackage{amsmath,amssymb,amsfonts,mathtools}
\allowdisplaybreaks
\setlength{\jot}{5pt}
\numberwithin{equation}{section}

\title{Three-$\rho$ finite terms with the complex conjugate added}
\author{}
\date{}

\begin{document}
\maketitle
\tableofcontents

\section{What adding the conjugate does}

Every block of the final result of \texttt{Three\_rho\_limiting.tex} has the form
\begin{align}
X_{b}=i\,\mathcal{P}_{0}\int g_{b}\;\mathcal{O}_{b},\qquad\mathcal{P}_{0}\equiv\frac{1}{(2\pi)^{3}}\,\frac{g^{4}}{8\pi^{5}}\,\frac{1}{k^{+}},
\end{align}
with a $c$-number kernel $g_{b}$ (phases included) and an operator $\mathcal{O}_{b}$ made of Wilson lines and $\rho$'s. The Wilson lines $U^{ab}$ and $f^{abc}$ are real and $\rho$ is hermitian, so the conjugate reverses the order of the $\rho$'s and conjugates $g_{b}$. With
\begin{align}
\mathsf{H}[\mathcal{O}]\equiv\tfrac12\big(\mathcal{O}+\mathcal{O}^{\dagger}\big),\qquad\mathsf{A}[\mathcal{O}]\equiv\tfrac12\big(\mathcal{O}-\mathcal{O}^{\dagger}\big),
\end{align}
\begin{align}
\boxed{\;X_{b}+X_{b}^{\dagger}=2\,\mathcal{P}_{0}\int\Big\{-{\rm Im}\,g_{b}\;\mathsf{H}[\mathcal{O}_{b}]+i\,{\rm Re}\,g_{b}\;\mathsf{A}[\mathcal{O}_{b}]\Big\}\;} \label{eq:rule}
\end{align}
So for each operator structure exactly half of the kernel survives: the hermitian part of the operator keeps only ${\rm Im}\,g$, the anti-hermitian part only ${\rm Re}\,g$. The other half cancels against the conjugate. Nothing else cancels, because the conjugate has the roles of amplitude and conjugate amplitude exchanged and cannot be relabelled into $X_{b}$.

\section{Hermitian and anti-hermitian parts of the operators}

For hermitian $\rho_{1},\rho_{2},\rho_{3}$ (any positions and colour indices) and a real coefficient:
\begin{align}
\mathsf{H}[\rho_{1}\rho_{2}\rho_{3}]&=\tfrac12\big(\rho_{1}\rho_{2}\rho_{3}+\rho_{3}\rho_{2}\rho_{1}\big), &
\mathsf{A}[\rho_{1}\rho_{2}\rho_{3}]&=\tfrac12\big(\rho_{1}[\rho_{2},\rho_{3}]+[\rho_{1},\rho_{3}]\rho_{2}+\rho_{3}[\rho_{1},\rho_{2}]\big),\notag\\
\mathsf{H}\big[\rho_{1}[\rho_{2},\rho_{3}]\big]&=\tfrac12\big[\rho_{1},[\rho_{2},\rho_{3}]\big], &
\mathsf{A}\big[\rho_{1}[\rho_{2},\rho_{3}]\big]&=\tfrac12\big\{\rho_{1},[\rho_{2},\rho_{3}]\big\},\notag\\
\mathsf{H}\big[[\rho_{1},\rho_{2}]\rho_{3}\big]&=\tfrac12\big[[\rho_{1},\rho_{2}],\rho_{3}\big], &
\mathsf{A}\big[[\rho_{1},\rho_{2}]\rho_{3}\big]&=\tfrac12\big\{[\rho_{1},\rho_{2}],\rho_{3}\big\},\notag\\
\mathsf{H}\big[\rho_{1}\{\rho_{2},\rho_{3}\}\big]&=\tfrac12\big\{\rho_{1},\{\rho_{2},\rho_{3}\}\big\}, &
\mathsf{A}\big[\rho_{1}\{\rho_{2},\rho_{3}\}\big]&=\tfrac12\big[\rho_{1},\{\rho_{2},\rho_{3}\}\big]. \label{eq:HA}
\end{align}
The operators of the final result (notation of the previous note) and the rule that applies:
\begin{center}
\renewcommand{\arraystretch}{1.35}
\begin{tabular}{l|l|l}
operator & $\rho$ content & rule in \eqref{eq:HA}\\ \hline
$\Omega$ (four terms) & $[\rho(x'),\rho(y)]\,\rho(x)$ and $\rho(x')\,[\rho(y),\rho(x)]$ & lines 2, 3\\
$\mathcal{Z}=\mathcal{D}_{0}^{a}f^{adc}U^{db}(x)\Delta U^{ce}(z,y)[\rho^{e}(y),\rho^{b}(x)]$ & $\rho(x')\,[\rho(y),\rho(x)]$ & line 2\\
$\mathcal{W}=\mathcal{D}_{0}^{a}f^{abd}U^{de}(x)\Delta U^{bc}(z,y)\rho^{e}(x)\rho^{c}(y)$ & $\rho(x')\rho(x)\rho(y)$ & line 1\\
$\mathcal{W}_{w}\equiv\mathcal{W}$ with $\Delta U^{bc}(z,y)\to\Delta U^{bc}(w,y)$ & $\rho(x')\rho(x)\rho(y)$ & line 1\\
$\mathcal{O}_{6}=\mathcal{D}_{0}^{a}f^{adc}\Delta U^{db}(x,y')\Delta U^{ce}(z,y)\{\rho^{e}(y),\rho^{b}(y')\}$ & $\rho(x')\,\{\rho(y),\rho(y')\}$ & line 4\\
$\mathcal{C}(q)$ & $\rho(x')\rho(x)\rho(y)$ and $\rho(x')\,[\rho(y),\rho(x)]$ & lines 1, 2
\end{tabular}
\end{center}
For the commutator operators ($\Omega$, $\mathcal{Z}$, second part of $\mathcal{C}$): $\mathsf{A}$ is the operator with the product of the commutator and the third $\rho$ replaced by half their anticommutator, and $\mathsf{H}$ is a double commutator.

\section{Real and imaginary parts of the kernels}

Phases and arguments:
\begin{align}
&\varphi\equiv k\cdot(w'-w),\qquad\theta\equiv k\cdot(w'-z),\qquad\theta_{x}\equiv k\cdot(w'-x),\qquad\sigma\equiv k\cdot(z-x),\notag\\
&\tau\equiv k\cdot(r-z),\qquad\varphi_{7}\equiv k\cdot(y'-w),\qquad s\equiv k\cdot(y'-z).
\end{align}
Real functions:
\begin{align}
{\rm Cin}(t)\equiv\gamma_{E}+\log|t|-{\rm Ci}(|t|)\ \ (\text{even}),\qquad{\rm Si}(t)\ \ (\text{odd}),\qquad{\rm si}(s)\equiv\tfrac{\pi}{2}\,{\rm sgn}(s)-{\rm Si}(s)\ \ (\text{odd}),
\end{align}
\begin{align}
{\rm Ein}(\pm it)={\rm Cin}(t)\pm i\,{\rm Si}(t),\qquad
J_{0}(s)=\underbrace{-\cos s\,{\rm Ci}(|s|)+\sin s\,{\rm si}(s)}_{j_{R}(s)}+i\,\underbrace{\big(-\sin s\,{\rm Ci}(|s|)-\cos s\,{\rm si}(s)\big)}_{j_{I}(s)},
\end{align}
\begin{align}
&\gamma_{E}+\log(is)=\gamma_{E}+\log|s|+\tfrac{i\pi}{2}{\rm sgn}(s),\qquad\frac{i}{s-i0}=-\pi\,\delta(s)+i\,{\rm P}\frac1s,\notag\\
&e^{-i\theta}\,\frac{1-e^{-i\sigma}}{i\sigma}=\frac{\sin\theta_{x}-\sin\theta}{\sigma}+i\,\frac{\cos\theta_{x}-\cos\theta}{\sigma}.
\end{align}
With $u=n(x'-w')$, $v=n(y-z)$ and the kernels $K^{ij}$, $\mathcal{T}^{ij}$, $L$, $L_{d}$ of the previous note:
\begin{center}
\renewcommand{\arraystretch}{1.6}
\begin{tabular}{c|l|l|l}
$b$ & $\mathcal{O}_{b}$ & ${\rm Re}\,g_{b}$ & ${\rm Im}\,g_{b}$\\ \hline
1 & $\Omega$ & $\cos\varphi\;u^{i}v^{j}K^{ij}L$ & $-\sin\varphi\;u^{i}v^{j}K^{ij}L$\\
2 & $\mathcal{Z}$ & $\tfrac12\cos\varphi\;u^{i}v^{j}\mathcal{T}^{ij}$ & $-\tfrac12\sin\varphi\;u^{i}v^{j}\mathcal{T}^{ij}$\\
3 & $\mathcal{W}-\tfrac12\mathcal{Z}$ & $-u\cdot v\displaystyle\int_{r}N_{r}\big[\cos\theta\,{\rm Cin}(\tau)-\sin\theta\,{\rm Si}(\tau)\big]$ & $u\cdot v\displaystyle\int_{r}N_{r}\big[\sin\theta\,{\rm Cin}(\tau)+\cos\theta\,{\rm Si}(\tau)\big]$\\
4 & $\mathcal{W}_{w}$ & $-\tfrac12\cos\varphi\;L_{d}\,N_{4}$ & $\tfrac12\sin\varphi\;L_{d}\,N_{4}$\\
5 & $\mathcal{W}$ & $\pi\sin\varphi\;u^{i}v^{j}K^{ij}$ & $\pi\cos\varphi\;u^{i}v^{j}K^{ij}$\\
6 & $\mathcal{O}_{6}$ & $\tfrac14\mathcal{K}_{1}\frac{\sin\theta_{x}-\sin\theta}{\sigma}+\tfrac12\mathcal{K}_{2}\big[\cos\theta\,{\rm Cin}(\sigma)+\sin\theta\,{\rm Si}(\sigma)\big]$ & $\tfrac14\mathcal{K}_{1}\frac{\cos\theta_{x}-\cos\theta}{\sigma}+\tfrac12\mathcal{K}_{2}\big[\cos\theta\,{\rm Si}(\sigma)-\sin\theta\,{\rm Cin}(\sigma)\big]$
\end{tabular}
\end{center}
\begin{align}
&N_{r}\equiv n(x-r)\cdot n(z-r),\qquad N_{4}\equiv u\cdot n(y-w)\;\,n(x-z)\cdot n(w-z),\notag\\
&\mathcal{K}_{1}\equiv u\cdot n(z-x)\;\,v\cdot n(y'-x),\qquad\mathcal{K}_{2}\equiv u\cdot v\;\,n(z-x)\cdot n(y'-x).
\end{align}
Block 5 is $-\frac{g^{4}f^{abd}}{8\pi^{4}}K\,\mathcal{W}$ written as $i\mathcal{P}_{0}\,g_{5}\,\mathcal{W}$ with $g_{5}=i\pi e^{-i\varphi}u^{i}v^{j}K^{ij}$. Block 7 has three kernel functions $F$ multiplying different operators; for each, ${\rm Re}(e^{-i\varphi_{7}}F)=\cos\varphi_{7}F_{R}+\sin\varphi_{7}F_{I}$ and ${\rm Im}(e^{-i\varphi_{7}}F)=\cos\varphi_{7}F_{I}-\sin\varphi_{7}F_{R}$:
\begin{center}
\renewcommand{\arraystretch}{1.5}
\begin{tabular}{l|l|c|c}
tensor structure & operator & $F_{R}$ & $F_{I}$\\ \hline
$M_{a}\equiv n(y'-z)\cdot n(x'-y')\;\,n(y-w)\cdot n(x-z)$ & $\mathcal{C}(y')$ & $j_{R}(s)$ & $j_{I}(s)$\\
$M_{b}\equiv n(y'-z)\cdot n(x-z)\;\,n(y-w)\cdot n(x'-y')$ & $\mathcal{C}(y')$ & $\gamma_{E}+\log|s|$ & $\tfrac{\pi}{2}{\rm sgn}(s)$\\
$M_{c}\equiv n(y'-z)\cdot n(y-w)$ & $\mathcal{C}_{\Delta}$ & $-\pi\delta(s)$ & ${\rm P}\frac1s$
\end{tabular}
\end{center}
with $\mathcal{C}_{\Delta}\equiv n(x-z)\cdot n(x'-y')\,\mathcal{C}(y')-n(x-z)\cdot n(x'-z)\,\mathcal{C}(z)$.

\section{Result}

Inserting the tables in \eqref{eq:rule}:
\begin{align}
\boxed{\;\begin{aligned}
\frac{d^{3}N^{3\rho}_{\rm NLO}}{d^{3}k}&=2\mathcal{P}_{0}\int_{x,y,z,w,x',w'}u^{i}v^{j}\Big\{K^{ij}L\Big(\sin\varphi\;\mathsf{H}[\Omega]+i\cos\varphi\;\mathsf{A}[\Omega]\Big)+\tfrac12\mathcal{T}^{ij}\Big(\sin\varphi\;\mathsf{H}[\mathcal{Z}]+i\cos\varphi\;\mathsf{A}[\mathcal{Z}]\Big)\Big\}\\
&-2\mathcal{P}_{0}\int_{x,y,z,x',w'}u\cdot v\int_{r}N_{r}\Big\{\big[\sin\theta\,{\rm Cin}(\tau)+\cos\theta\,{\rm Si}(\tau)\big]\mathsf{H}\big[\mathcal{W}-\tfrac12\mathcal{Z}\big]\\
&\hspace{4.2cm}+i\big[\cos\theta\,{\rm Cin}(\tau)-\sin\theta\,{\rm Si}(\tau)\big]\mathsf{A}\big[\mathcal{W}-\tfrac12\mathcal{Z}\big]\Big\}\\
&-\mathcal{P}_{0}\int_{x,y,z,w,x',w'}L_{d}\,N_{4}\Big\{\sin\varphi\;\mathsf{H}[\mathcal{W}_{w}]+i\cos\varphi\;\mathsf{A}[\mathcal{W}_{w}]\Big\}\\
&-2\pi\mathcal{P}_{0}\int_{x,y,z,w,x',w'}u^{i}v^{j}K^{ij}\Big\{\cos\varphi\;\mathsf{H}[\mathcal{W}]-i\sin\varphi\;\mathsf{A}[\mathcal{W}]\Big\}\\
&+2\mathcal{P}_{0}\int_{x,z,y,x',y',w'}\Big\{-{\rm Im}\,g_{6}\;\mathsf{H}[\mathcal{O}_{6}]+i\,{\rm Re}\,g_{6}\;\mathsf{A}[\mathcal{O}_{6}]\Big\}\\
&+2\mathcal{P}_{0}\int_{w,x',y',z,y,x}\sum_{F}M_{F}\Big\{\big[\sin\varphi_{7}F_{R}-\cos\varphi_{7}F_{I}\big]\mathsf{H}[\mathcal{C}_{F}]\\
&\hspace{4.2cm}+i\big[\cos\varphi_{7}F_{R}+\sin\varphi_{7}F_{I}\big]\mathsf{A}[\mathcal{C}_{F}]\Big\}
\end{aligned}\;} \label{eq:result}
\end{align}
(${\rm Re}\,g_{6}$, ${\rm Im}\,g_{6}$ from the table; in the last line $F$ runs over the three rows of the block-7 table and $\mathcal{C}_{F}$ is the operator of that row.) This expression is hermitian term by term: every $\mathsf{H}$ is hermitian with a real coefficient and every $i\,\mathsf{A}$ is hermitian with a real coefficient. There is no ``${\rm c.c.}$'' left to add.

\section{What cancelled}

\begin{itemize}
\item \textbf{Blocks 1 and 2} (commutators). $\mathsf{H}[\Omega]$ and $\mathsf{H}[\mathcal{Z}]$ are double commutators. Up to those, the $\sin[k\cdot(w'-w)]$ half cancels and only $i\cos\varphi\,\mathsf{A}$ remains.
\item \textbf{Blocks 3, 4, 6, 7} ($i\times$ three $\rho$). For the symmetrised three-$\rho$ operator $\mathsf{H}$ the real part of the kernel cancels and only ${\rm Im}\,g$ remains; for block 4 this is $\sin\varphi$.
\item \textbf{Block 5} (no $i$). The opposite happens: the $\sin\varphi$ half cancels for $\mathsf{H}[\mathcal{W}]$ and $\cos\varphi$ remains.
\item The commutator pieces inside blocks 3 and 7 behave like blocks 1 and 2.
\end{itemize}

\section{Sorting by the number of $\rho$'s}

With $[\rho^{a}(x),\rho^{b}(y)]=i\kappa f^{abc}\rho^{c}(x)\delta^{2}(x-y)$ every commutator lowers the number of $\rho$'s by one and supplies a factor $i\kappa$.
\begin{center}
\renewcommand{\arraystretch}{1.4}
\begin{tabular}{l|l|l}
order & origin in \eqref{eq:result} & coefficient\\ \hline
three $\rho$, symmetrised & $\mathsf{H}$ of $\mathcal{W}$, $\mathcal{W}_{w}$, $\mathcal{O}_{6}$ and of the first part of $\mathcal{C}$ & $-2\mathcal{P}_{0}\,{\rm Im}\,g$\\
two $\rho$ & $i\,\mathsf{A}$ of $\Omega$, $\mathcal{Z}$, second part of $\mathcal{C}$ & $-2\kappa\mathcal{P}_{0}\,{\rm Re}\,g\times\tfrac12\{\rho,\rho\}$\\
two $\rho$ & $i\,\mathsf{A}$ of $\mathcal{W}$, $\mathcal{W}_{w}$, $\mathcal{O}_{6}$, first part of $\mathcal{C}$ & $-\kappa\mathcal{P}_{0}\,{\rm Re}\,g\times(\text{ordered }\rho\rho)$\\
one $\rho$ & $\mathsf{H}$ of $\Omega$, $\mathcal{Z}$, second part of $\mathcal{C}$ (double commutators) & $+\kappa^{2}\mathcal{P}_{0}\,{\rm Im}\,g$
\end{tabular}
\end{center}
Explicitly, for the two rules used most:
\begin{align}
i\,\mathsf{A}\big[\rho^{a}(x')\,[\rho^{b}(y),\rho^{c}(x)]\big]&=-\frac{\kappa}{2}\,f^{bcg}\,\big\{\rho^{a}(x'),\rho^{g}(x)\big\}\,\delta^{2}(x-y),\notag\\
i\,\mathsf{A}\big[\rho^{a}(x')\rho^{b}(x)\rho^{c}(y)\big]&=-\frac{\kappa}{2}\Big(f^{bcg}\rho^{a}(x')\rho^{g}(x)\,\delta^{2}(x-y)+f^{acg}\rho^{g}(x')\rho^{b}(x)\,\delta^{2}(x'-y)+f^{abg}\rho^{c}(y)\rho^{g}(x')\,\delta^{2}(x'-x)\Big).
\end{align}
For blocks 1 and 2 the colour reduction of the previous note applies with the two $\rho$'s symmetrised, e.g.
\begin{align}
i\,\mathsf{A}[\mathcal{Z}]=-\frac{\kappa}{2}\,\delta^{2}(x-y)\,\Delta U^{ab'}(x',w')\Big[-f^{adc}f^{gbe}U^{db}(x)U^{ce}(z)+N_{c}U^{ag}(x)\Big]\big\{\rho^{b'}(x'),\rho^{g}(x)\big\}.
\end{align}

\section{Status of the three-$\rho$ part}

After adding the conjugate, the genuinely three-$\rho$ part is $-2\mathcal{P}_{0}\,{\rm Im}\,g_{b}\,\mathsf{H}[\mathcal{O}_{b}]$ for $b=3$ ($\mathcal{W}$ piece), 4, 5, 6, 7 (first piece).
\begin{itemize}
\item For $b=3,4,6,7$ the surviving ${\rm Im}\,g_{b}$ is odd under a global reflection of all coordinates ($\sin\varphi$; $\sin\theta\,{\rm Cin}+\cos\theta\,{\rm Si}$; $\cos\varphi_{7}F_{I}-\sin\varphi_{7}F_{R}$ with $F_{I}$ odd and $F_{R}$ even; \dots). These pieces vanish after the transverse integrals if the averages over projectile and target are reflection symmetric. That is an assumption about the averages; there is no cancellation point by point.
\item For $b=5$ the surviving kernel is $\pi\cos\varphi\,K^{ij}$, which is even. This term does not vanish by symmetry.
\item The large-distance divergences noted before (the $r$ integral of block 3 with block 4, and the $\gamma_{E}+\log|s|$ piece of block 7) are not affected by adding the conjugate.
\end{itemize}

\section*{Checks}
The four rules \eqref{eq:HA} and \eqref{eq:rule} were checked with random hermitian matrices ($10^{-14}$); the real and imaginary parts of ${\rm Ein}(\pm it)$, $J_{0}(s)$, $\gamma_{E}+\log(is)$ and $e^{-i\theta}h_{0}(\sigma)$ were checked numerically to 10 digits.

\end{document}
